% CLAUDE REPORT::WR-B3-CONSERVATION-0001::2026-11-14T13:50Z
% !TEX program = xelatex
% Compiles with XeLaTeX (UTF-8 source).
\documentclass[11pt]{article}
\usepackage[margin=1in]{geometry}
\usepackage{amsmath,amssymb,amsthm}
\usepackage[hidelinks]{hyperref}
\newtheorem{theorem}{Theorem}
\newtheorem{corollary}[theorem]{Corollary}
\newtheorem{lemma}[theorem]{Lemma}
\theoremstyle{definition}
\newtheorem{definition}[theorem]{Definition}
\newcommand{\Wbar}{\overline W}
\title{Where the difficulty of the worst-case derivation program sits:\\
B3 is equivalent to $\mathrm{SAT}\notin\mathrm{P}$ for every calculus with a known lower bound}
\author{Claude scheduled run WR-B3-CONSERVATION-0001 (14 November 2026, 05:50 PST)}
\date{}
\begin{document}
\maketitle

\noindent\textbf{Status summary.}
Theorem~\ref{thm:conservation} and Corollary~\ref{cor:weak} are \textbf{PROVED/EXACT} (elementary).
Theorem~\ref{thm:pv} is a \textbf{KNOWN THEOREM} (Cook 1975; Krajíček--Pudlák 1989), restated for $\Wbar_R$.
The unconditional bridge B3 for any calculus whose lower bound is still open is \textbf{OPEN}.
A superpolynomial lower bound on $\Wbar_R$ for any $R$ as strong as Extended Resolution is \textbf{OPEN}.
\textbf{P versus NP remains OPEN.}
No claim here strengthens any result in root \texttt{main.tex}.

\section{Target, as fixed by the shared directive}
We use the shared definitions from
\path{research/coordination/2026-11-02_1204_WORST_CASE_MINIMUM_DERIVATION_FIVE_ROLE_REFOCUS.md}.
Fix a sound and complete refutation system $R$ for CNF.
Each derivation $\pi$ is a DAG with nonnegative node weights, and $\mathrm{cost}_R(\pi)$ is the total weight of its distinct nodes.
Then
\[
W_R(F)=\min_{\pi:F\vdash_R\bot}\mathrm{cost}_R(\pi),\qquad
\Wbar_R(N)=\max_{F\in\mathrm{UNSAT},\ |\langle F\rangle|\le N}W_R(F).
\]
The inner minimum grants free, perfect choice of the best reduction sequence, as the user's directive requires.
In the multisign DAG-hypergraph picture, it is the cheapest legal sequence that drives the hypergraph to the minimal contradictory reduced state.
The proof-search cost is never charged.
Write $\mathrm{LB}_R$ for the statement that $\Wbar_R$ is superpolynomial, and
\[
\mathrm{B3}_R:\qquad \mathrm{SAT}\in\mathrm{P}\ \Longrightarrow\ \exists C,k\ \forall N:\ \Wbar_R(N)\le C(N+1)^k .
\]

\section{Main observation}
\begin{theorem}[B3 conservation]\label{thm:conservation}
For every fixed system $R$, whatever its rules and weights:
\begin{enumerate}
\item $\mathrm{SAT}\notin\mathrm{P}\Rightarrow \mathrm{B3}_R$.
\item $\mathrm{LB}_R\wedge\mathrm{B3}_R\Rightarrow\mathrm{SAT}\notin\mathrm{P}$.
\item Consequently, if $\mathrm{LB}_R$ is an established theorem, then $\mathrm{B3}_R\Longleftrightarrow \mathrm{SAT}\notin\mathrm{P}$.
\end{enumerate}
\end{theorem}
\begin{proof}
(1) If the antecedent of $\mathrm{B3}_R$ is false, the implication holds vacuously.
(2) This is the contrapositive of $\mathrm{B3}_R$: a superpolynomial $\Wbar_R$ contradicts its consequent.
(3) Combine (1) and (2).
\end{proof}

\begin{corollary}[weak calculi carry no independent bridge]\label{cor:weak}
Let $R$ be any of the following, with unit node weights and polynomial-size encodings:
tree-like or DAG Resolution (Haken 1985: pigeonhole);
Res$(k)$ for constant $k$;
Cutting Planes (Pudlák 1997: clique--colouring);
bounded-depth Frege (Ajtai 1988; Beame--Impagliazzo--Pitassi, Krajíček--Pudlák--Woods 1995: pigeonhole);
Polynomial Calculus over a fixed field (Razborov 1998; Buss--Grigoriev--Impagliazzo--Pitassi 1999: Tseitin mod $q$).
Also include any fragment of the corpus calculus $R_{\mathrm{mix341}}$ that is polynomially simulated by one of these systems.
For each such $R$, $\mathrm{LB}_R$ is a known theorem.
Therefore proving $\mathrm{B3}_R$ is \emph{exactly as hard as proving} $\mathrm{SAT}\notin\mathrm{P}$.
\end{corollary}
\begin{proof}
Theorem~\ref{thm:conservation}(3), applied to the cited lower bounds.
The encoded size $N$ is polynomial in the number of clauses of each cited family, so the lower bounds, which are exponential in the variable count, stay superpolynomial in $N$.
\end{proof}

\paragraph{Interpretation.}
B3 is the logically weakest bridge, but it does not divide the difficulty of the problem.
Splitting $\mathrm{SAT}\notin\mathrm{P}$ into ``$\mathrm{LB}_R$'' plus ``$\mathrm{B3}_R$'' is a real reduction only if both of the following hold:
\begin{itemize}
\item $\mathrm{LB}_R$ is open, so $R$ is at least as strong as every system with a known superpolynomial lower bound.
\item $\mathrm{B3}_R$ is provable without first proving $\mathrm{SAT}\notin\mathrm{P}$, so the argument must use the hypothesis $\mathrm{SAT}\in\mathrm{P}$ constructively.
\end{itemize}
The second condition means a valid proof of $\mathrm{B3}_R$ must take a hypothetical polynomial-time SAT algorithm $A$ and build short $R$-refutations from it.
It cannot use a lower-bound argument on $R$.
For weak $R$, any such construction would contradict the known lower bounds.
So for weak $R$ the only true proofs of $\mathrm{B3}_R$ are proofs of $\mathrm{SAT}\notin\mathrm{P}$ in disguise.
This sharpens the corpus's existing B1 refutation, the parity recognizer against Resolution.
For weak calculi, even B3 cannot be established by any ``translation'' argument.

\section{Strong calculi: the remaining gap is truth versus feasible provability}
The corpus calculus $R_{\mathrm{mix341}}$ contains RES and VARDEF, that is, Extended Resolution, which is p-equivalent to Extended Frege (EF).
No superpolynomial lower bound is known for EF, so $\mathrm{LB}_{\mathrm{mix341}}$ is open.
This is the regime where the split could be a real reduction.

\begin{theorem}[KNOWN THEOREM: Cook 1975; Krajíček--Pudlák 1989; see Krajíček 1995]\label{thm:pv}
Let $A$ be a polynomial-time SAT decider.
Let $\mathrm{Ref}_{A,n}$ be the polynomial-size propositional formulas stating that $A$ never rejects a satisfiable input of length $n$.
Suppose Cook's theory $\mathrm{PV}$ proves the $\forall\Pi^b_1$ correctness sentence of $A$.
Then $\mathrm{Ref}_{A,n}$ has EF proofs of size $n^{O(1)}$.
Hence EF polynomially simulates the proof system ``$\pi$ certifies $F$ when $A$ rejects $F$''.
Hence there are constants with $\Wbar_{\mathrm{EF}}(N)\le C(N+1)^k$, and the same holds for any $R$ that polynomially simulates EF, including $R_{\mathrm{mix341}}$ under the corpus encodings.
\end{theorem}
\noindent This is the standard propositional translation of $\mathrm{PV}$ into EF together with the reflection simulation argument.
We restate it; we do not claim it as new.

\begin{corollary}[split for EF-strength $R$]\label{cor:split}
For $R\supseteq\mathrm{ER}$:
\[
\mathrm{B3}^{\mathrm{PV}}_R:\quad\bigl(\exists A\ \text{polytime},\ \mathrm{PV}\vdash \mathrm{Correct}(A)\bigr)\Longrightarrow \Wbar_R\ \text{polynomial}
\]
is a \textbf{KNOWN THEOREM}.
Therefore $\mathrm{LB}_R$ alone already proves that no polynomial-time SAT algorithm is \emph{PV-provably} correct.
The unconditional $\mathrm{B3}_R$ asks for more.
It requires short $R$-proofs of $\mathrm{Ref}_{A,n}$ for an algorithm that is \emph{correct} but whose correctness need not be provable in any fixed feasible theory.
That gap between true and feasibly provable correctness is \textbf{OPEN}.
It is the same gap the corpus isolated as LSPV-1/337 (finite reflection in a fixed theory $T$) and LSPV-1/353--354 (single reflection family $H_{A,\ell}$).
\end{corollary}

\section{Consequences for the five assignments}
\begin{itemize}
\item \textbf{General SAT Lower-Bound Lead:} do not present B3 for Resolution, CP, PC or AC$^0$-Frege fragments as a ``weaker'' target.
By Corollary~\ref{cor:weak} it is equivalent to the full separation.
The only productive B3 work is for EF-strength $R$, and it must close the gap between true and feasibly provable correctness.
\item \textbf{SAT-Circuit Reduction Builder:} any certified shortest-DAG hardness for a family is a family-level fact about $R$.
It contributes to $\mathrm{LB}_R$ only if the family remains hard after VARDEF/ER is added.
Pigeonhole and Tseitin formulas are hard for Resolution but have polynomial-size EF refutations (Cook--Reckhow 1979 for PHP; Gaussian elimination for Tseitin), so they are not candidates for $\mathrm{LB}_R$.
Random 3-CNF is not known to be easy in EF, and it is not known to be hard either.
\item \textbf{Hypergraph \& Proof Certificates:} any family proposed as hard for EF-strength $R$ needs, at minimum, a candidate that is not known to be easy in EF.
Random 3-CNF and the circuit-lower-bound formulas of Razborov and Krajíček are the standard conjectured candidates.
Any proof would be a major new EF lower bound and must therefore get past the known meta-barriers (for example the Pich--Santhanam unprovability line).
\item \textbf{Time-Space Resource Auditor:} Theorem~\ref{thm:pv} charges only proof size.
The $\mathrm{PV}\to$EF translation size is polynomial with exponent depending on the $\mathrm{PV}$ proof, which is fixed per $A$, so B3's ``$C,k$ may depend on $A$'' quantifier order is respected.
\item \textbf{Independent Red Team:} treat any claimed proof of $\mathrm{B3}_R$ for an $R$ with a known lower bound as a claimed proof of P$\neq$NP.
Audit it at that level, including relativization: Theorem~\ref{thm:conservation} relativizes, so any argument proving $\mathrm{B3}_R$ for such $R$ must be nonrelativizing.
\end{itemize}

\section*{CORPUS REVIEWED}
\begin{itemize}
\item \texttt{AGENTS.md} (all allocation headings, including the 2026-11-09 knot/link handoff note).
\item The shared directive \path{research/coordination/2026-11-02_1204_WORST_CASE_MINIMUM_DERIVATION_FIVE_ROLE_REFOCUS.md}.
\item The newest B3 synthesis \path{research/coordination/2026-11-13_1148_LSPV5_D1_B3_FULL_CHAIN_FIVE_ROLE_HANDOFF527.md}.
\item The TMR-5 handoff (014--016), and the TMR-1..5 commit stream on 2026-11-14 through \texttt{6fb855fd}.
\item Root \texttt{main.tex} regions LSPV1 FINITE-REFLECTION-BRIDGE337, UNIVERSAL-REFLECTION-SINGLE-FAMILY353 and PIN-AND-REPLAY354.
\item A repository-wide search for earlier ``B3 equivalent'' and ``vacuous'' formulations; the hits were LSPV-1/444 and handoffs 380, 417 and 419--422.
\item The earlier Claude notes in \path{research/agents/claude/}.
\end{itemize}

\section*{DUPLICATION / CONTRADICTION AUDIT}
The search found vacuity remarks about witness maps on negative inputs (LSPV-1/444).
It found no explicit statement of Theorem~\ref{thm:conservation}(3) or of Corollary~\ref{cor:weak}.
The reflection material in Theorem~\ref{thm:pv} overlaps LSPV-1/337 and LSPV-1/353--354.
It is cited here as a known theorem and is not claimed as new.
No contradiction with any verified corpus claim was found.

\section*{WHAT CHANGED RELATIVE TO ROOT MAIN}
Nothing in root \texttt{main.tex} was edited.
This run is a Claude reference note outside the LSPV/TMR stagger and owns no \texttt{main.tex} region.
For LSPV-5/TMR-5, the suggested synthesis is: label B3 for every calculus with a known lower bound as ``equivalent to $\mathrm{SAT}\notin\mathrm{P}$; not a reduction of difficulty''.

\section*{ADOPTED / REJECTED / RETIRED PRIOR COMPONENTS}
\begin{itemize}
\item \textbf{Adopted:} the directive's definitions; the LSPV-1/337 and /353 reflection framing.
\item \textbf{Rejected as framing:} ``B3 is the easy half''. It is the easy half only logically, not in difficulty.
\item \textbf{Retired:} nothing.
\end{itemize}

\section*{STALE OR OVERSTATED CLAIMS FOUND}
None was found in the files reviewed.

\section*{Barrier-transfer and literature preflight}
\begin{itemize}
\item \textbf{Obstacles.} (a) No EF lower bound. (b) The gap between true and feasibly provable correctness in B3.
\item \textbf{Directly transferable known theorems.} The Cook $\mathrm{PV}\to$EF translation and the reflection simulation (Theorem~\ref{thm:pv}).
\item \textbf{Structural analogy needing a new transfer theorem.} The Pich--Santhanam program (ECCC TR23-199) proves that EF lower bounds for EF augmented with witnessing axioms $w^k_n(f)$ imply circuit lower bounds for SAT, conditional on those formulas being tautologies.
That is a B3-type bridge with an added axiom scheme.
Transferring it to $R_{\mathrm{mix341}}$ would require the scheme to be admissible in the fixed calculus, which is not established.
\item \textbf{Failed transfer.} Theorem~\ref{thm:conservation} shows that no pure lower-bound technique for a weak calculus can supply B3.
\item \textbf{Post-quantum ledger.} Items 1--4 are not relevant to this meta-level statement.
Item 6 (collision extraction) is an analogy only: a binding-style extraction of a SAT witness from a rejected-but-satisfiable input is exactly what $\mathrm{Ref}_{A,n}$ asserts, and it gives no unconditional theorem.
\end{itemize}

\section*{Literature checked this run}
\begin{itemize}
\item S.~A.~Cook, Feasibly constructive proofs and the propositional calculus, STOC 1975.
\item J.~Krajíček and P.~Pudlák, Propositional proof systems, the consistency of first order theories and the complexity of computations, J.~Symbolic Logic 54 (1989).
\item J.~Krajíček, \emph{Bounded Arithmetic, Propositional Logic and Complexity Theory}, CUP 1995.
\item A.~Haken, The intractability of resolution, TCS 39 (1985).
\item P.~Pudlák, Lower bounds for resolution and cutting plane proofs and monotone computations, JSL 62 (1997).
\item A.~Razborov, Lower bounds for the polynomial calculus, Comput.\ Complexity 7 (1998).
\item J.~Pich and R.~Santhanam, Strong co-nondeterministic lower bounds for NP cannot be proved feasibly, STOC 2021.
\item J.~Pich and R.~Santhanam, Towards P$\neq$NP from Extended Frege lower bounds, ECCC TR23-199 (2023).
\item J.~Li, ECCC TR25-187 (2025), on PV provability and propositional translations. Only the abstract was consulted; the details are not relied on.
\end{itemize}

\section*{Process notes}
\begin{itemize}
\item No experiments were run, so the Python runtime rule did not apply.
\item The predecessor write-recovery check does not apply: this Claude routine is not part of the five-role stagger.
\item The five TMR roles' 2026-11-14 commits were all present at preflight.
\item Exactly five roles, their schedules and ownership were left unchanged, and no tasks were mutated.
\end{itemize}

\end{document}
% END_OF_VERIFIED_REPORT::CLAUDE::WR-B3-CONSERVATION-0001
